\section{Gaussian order derivatives and finite correction formulas}
\label{sec:kernels}

For $\Rea a>-1$ and real $\xi$, define
\begin{equation}\label{eq:H-definition}
 \HH_a(\xi)=\frac1{2\pi}\int_{\R}
     e^{-i\xi y-y^2/2}(-iy)^a\,dy,
 \qquad \HH_a^{[j]}(\xi)=\partial_a^j\HH_a(\xi).
\end{equation}
The logarithm convention is \eqref{eq:fourier-log}. All fixed order
derivatives are legitimate by domination at zero and infinity. For
real $a$ and $\xi$ the two halves of the integral are conjugates.

\begin{proposition}[Parabolic cylinder and Gamma forms]
\label{prop:H}
In the standard parabolic-cylinder normalization,
\begin{align}
 \HH_a(\xi)&=\frac{e^{-\xi^2/4}}{\sqrt{2\pi}}D_a(-\xi),
       \label{eq:H-cylinder}\\
 \HH_a(0)&=\frac{2^{(a-1)/2}}{\Gamma((1-a)/2)},
 &\partial_\xi\HH_a&=\HH_{a+1}.
       \label{eq:H-zero}
\end{align}
In particular $\HH_0=\normal$ and
$\HH_{2k}=\normal^{(2k)}$. Each function
$\HH_a^{[j]}$ used below is an ordinary convergent integral, even
when a later asymptotic series containing it does not converge.
\end{proposition}

\begin{proof}
The derivative identity follows immediately under the integral.
Splitting into $y>0$ and $y<0$, then using the Gamma integral,
gives
\[
 \HH_a(0)=\frac{2^{(a-1)/2}}{\pi}
       \Gamma((a+1)/2)\cos(\pi a/2).
\]
The reflection formula gives \eqref{eq:H-zero}. Integration by
parts in $y$, with vanishing boundary term
$y(-iy)^ae^{-y^2/2-i\xi y}$, gives
\[
 \HH_a''+\xi\HH_a'+(a+1)\HH_a=0.
\]
The function on the right of \eqref{eq:H-cylinder} satisfies the same
equation by the parabolic-cylinder differential equation.
The values of $D_a(0)$ and $D_a'(0)$
\cite[12.2.4--12.2.7]{dlmf-cylinder} give the same value and first
derivative at $\xi=0$, namely
$2^{(a-1)/2}/\Gamma((1-a)/2)$ and
$2^{a/2}/\Gamma(-a/2)$. Uniqueness for the ordinary differential
equation proves the identity.
\end{proof}

\subsection{The correction polynomials in two small parameters}

Define the finite functions
\begin{align}
 \mathcal C_0(\xi;q,\eps)&=\normal(\xi),\notag\\
 \mathcal C_j(\xi;q,\eps)
 &=\frac1{j!}\sum_{k=1}^j
     \stir{j}{k}\left(-\frac q2\right)^k\eps^{j-k}
           \HH_{2k}^{[j]}(\xi),\qquad j\ge1.
           \label{eq:correction}
\end{align}
They are homogeneous of degree $j$ in $(q,\eps)$. The first three are
\begin{align}
 \mathcal C_1&=-\frac q2\HH_2^{[1]},\label{eq:C1}\\
 \mathcal C_2&=-\frac{q\eps}{4}\HH_2^{[2]}
                  +\frac{q^2}{8}\HH_4^{[2]},\label{eq:C2}\\
 \mathcal C_3&=-\frac{q\eps^2}{12}\HH_2^{[3]}
             +\frac{q^2\eps}{8}\HH_4^{[3]}
             -\frac{q^3}{48}\HH_6^{[3]}.
             \label{eq:C3}
\end{align}
The spatial argument $\xi$ is suppressed on the right. These expressions
involve order derivatives of the Gaussian kernel; in general they are
not a Gaussian times an ordinary polynomial in $\xi$.

There is also a finite divided-difference form. For a holomorphic
function of $a$, put
\begin{equation}\label{eq:divided-difference}
 \mathfrak D_\eps^k F(a)
 =\frac1{\eps^k}\sum_{h=0}^k(-1)^{k-h}\binom{k}{h}F(a+h\eps),
 \qquad \mathfrak D_0^k F=F^{(k)}.
\end{equation}
The apparent poles are removable. The identity
\begin{equation}\label{eq:stirling-generating}
 \left(\frac{e^{\eps z}-1}{\eps}\right)^k
 =k!\sum_{j\ge k}\stir{j}{k}\frac{\eps^{j-k}z^j}{j!}
\end{equation}
explains the appearance of the Stirling numbers. We will prove that
both the $q$-ordered finite expression
\begin{equation}\label{eq:finite-difference-approx}
 \sum_{k=0}^N\frac{(-q/2)^k}{k!}
       \left.\mathfrak D_\eps^k\HH_a(\xi)\right|_{a=2k}
\end{equation}
and the total-degree expression
$\sum_{j=0}^N\mathcal C_j(\xi;q,\eps)$ approximate the same lattice
law with error $O(L^{-N-1})$. They differ from one another only at
the omitted orders. Formula \eqref{eq:finite-difference-approx}
can be useful for computation away from $\eps=0$; near zero it is
better evaluated using derivatives or the removable integral formula.

\subsection{Finite Gamma moments at the center}

For integers $j,k\ge1$ set
\begin{equation}\label{eq:M-jk}
 M_j(k)=\Rea\left[
   (\partial_z+\log2+i\pi)^j\Gamma(z)
                      \right]_{z=k+1/2}.
\end{equation}
This is also the convergent moment
\begin{equation}\label{eq:M-integral}
 M_j(k)=\Rea\int_0^\infty e^{-t}t^{k-1/2}
                          (\log t+\log2+i\pi)^j\,dt.
\end{equation}

\begin{corollary}[Central coefficient calculus]
\label{cor:central-corrections}
For every $j\ge1$,
\begin{equation}\label{eq:central-C}
 \frac{\mathcal C_j(0;q,\eps)}{\normal(0)}
 =\frac1{2^j j!\sqrt\pi}
     \sum_{k=1}^j\stir{j}{k}q^k\eps^{j-k}M_j(k).
\end{equation}
In particular
\begin{equation}\label{eq:central-first}
 \frac{\mathcal C_1(0;q,\eps)}{\normal(0)}
 =\frac q4(2-\gamma-\log2).
\end{equation}
\end{corollary}

\begin{proof}
Differentiate the Gamma--cosine expression in the proof of
\cref{prop:H}. At $a=2k$ it gives
\[
 \HH_{2k}^{[j]}(0)=
       \frac{(-1)^k2^{k-j-1/2}}{\pi}M_j(k).
\]
Substitution in \eqref{eq:correction} cancels the two factors
$(-1)^k$ and proves \eqref{eq:central-C}. For $j=k=1$, use
$\Gamma(3/2)=\sqrt\pi/2$ and
$\psi(3/2)=2-\gamma-2\log2$.
\end{proof}

These formulas can be converted into polygamma values, harmonic numbers,
zeta values, and the Euler constant by repeated differentiation of
$\Gamma$ and the half-integer polygamma recurrences. The moment form
keeps all signs and factors of two in a single finite expression.
